\documentclass[11pt]{article}
\usepackage[a4paper,margin=25mm]{geometry}
\usepackage{amsmath,amssymb,amsthm,booktabs}
\newtheorem{theorem}{Theorem}
\newtheorem{lemma}{Lemma}
\newcommand{\Aut}{\operatorname{Aut}}
\newcommand{\HS}{\mathrm{HS}}
\title{Incidence invariants for the Higman--Sims geometries}
\author{L091 classification component}
\date{}
\begin{document}
\maketitle
Let $G=\HS$ in the acquired permutation representation of degree $100$.
Its order is $44\,352\,000$. We use the complete subgroup and panel
inventory and the separately reviewed enumeration of the firm, residually
connected, flag-transitive, residually weakly primitive, locally
two-transitive IP2 coset geometries for this action. In ranks at least two,
that enumeration gives $130$ families and $86$ classes under $\Aut(G)$
and permutations of the types. The rank counts are respectively
$(6,44,67,12,1)$ and $(6,30,41,8,1)$ for ranks two through six.
There are no higher ranks. The acting automorphism group is represented
faithfully by the acquired normalizer of order $88\,704\,000$.

The enumeration and its completeness are inputs here. In particular, the
present argument must still exclude an incidence isomorphism which does
not preserve the specified action of $G$. We refer to the numbered
representatives in the fixed, acquired file \texttt{full\_fusion03/result.g}.
These numbers are internal identifiers, not the source's diagram labels.

\begin{lemma}\label{shadows}
Suppose a group $G$ acts transitively on each type of a finite typed
incidence structure. Let $X$ be a type and let $S_y\subseteq X$ be the
set of objects incident with an object $y$ of another type. If
$G_{S_y}=G_y$, the map $gy\mapsto gS_y$ is injective. More generally,
suppose the objects of each remaining type can successively be
identified by their sets of neighbours in types already identified.
Then restriction to $X$ is faithful on the group of automorphisms
preserving these types.
\end{lemma}
\begin{proof}
The fibres of the equivariant map $G/G_y\longrightarrow G/G_{S_y}$
have size $[G_{S_y}:G_y]$. Equality of the stabilizers proves the first
claim. An automorphism fixing $X$ pointwise fixes every distinct subset
of $X$, hence fixes the next type pointwise. Induction proves the
second assertion. This argument does not require that $G$ be the full
automorphism group.
\end{proof}

\begin{lemma}\label{characteristic}
Let $\Gamma$ and $\Gamma'$ have a distinguished type $X$ and $X'$,
recognizable from their type sizes. Suppose each structure defines a
graph on that type purely by incidence, restriction of its automorphism
group to that type is faithful, and the image of its specified group
$G$ is the derived subgroup of the full graph automorphism group.
Every incidence isomorphism $\Gamma\longrightarrow\Gamma'$ conjugates
the specified $G$ actions.
\end{lemma}
\begin{proof}
An incidence isomorphism $f$ induces a graph isomorphism $f_X$. It
conjugates the full graph automorphism groups and consequently their
derived subgroups. For $g\in G$, choose $h\in G$ whose action on $X'$
equals that of $fgf^{-1}$. The automorphism $h^{-1}fgf^{-1}$ is trivial
on $X'$, so faithfulness makes it trivial on all of $\Gamma'$.
This gives the asserted conjugation of the entire actions.
\end{proof}

\begin{theorem}\label{classification}
Under the acquired enumeration and the exact computational inputs
specified below, the $86$ representatives are pairwise nonisomorphic
as incidence geometries, even allowing permutations of types. There
is one additional isomorphism class in rank one, consisting of a set
of $176$ objects. It is a truncation of a higher rank geometry.
Thus the numbers of isomorphism classes in ranks one through six are
\[
                  1,\quad6,\quad30,\quad41,\quad8,\quad1.
\]
\end{theorem}
\begin{proof}
Write $G_i$ for the stabilizers of a base chamber. Since the action is
flag-transitive, the number of flags whose set of types is a nonempty
subset $J$ is $|G|/|\bigcap_{j\in J}G_j|$. These counts are incidence
invariants, with their coordinates permuted when the types are
permuted. In particular the multiset of orders $|G_i|$, the multiset
of orders $|G_i\cap G_j|$, and the Borel order are invariants for
these fixed-order actions. Computing these from the actual subgroups
separates all representatives except the pairs $(14,15)$ and $(83,84)$.
For the rank five pair $(83,84)$, the vector of all nonempty flag
stabilizer orders, minimized over every permutation of the five types,
has different values. These are checked by intersections of the
acquired subgroups, not inferred from the printed diagrams.

It remains to separate $14$ and $15$. Their three stabilizer orders
are $192,1920,40320$. Their type sizes are therefore
$231000,23100,1100$, so every isomorphism preserves each type.
Call these types lines, blocks and points, respectively, and let
$P,Q,R$ denote their stabilizers at a base chamber. Take the actual
action on the $1100$ right cosets of $R$. It is faithful. If $z=R$,
the shadows of the base line and block on the point type are
$z^P$ and $z^Q$, of sizes two and sixteen.

Define two points to be adjacent when some line is incident with both.
This graph is intrinsic to the geometry. To construct it, let $w$ be
the other member of $z^P$. Its neighbours at $z$ are exactly $w^R$:
the lines incident with $z$ form the $R$ orbit of the base line.
Transporting this neighbour set by a representative of every coset
gives the whole graph. In both cases it has $1100$ vertices and
valency $105$.

The stabilizer in $G$ of the base block shadow $z^Q$ has order $1920$,
equal to $|Q|$. Thus blocks are determined by their sixteen point
shadows. Lines are \emph{not} determined by their two point shadows:
the corresponding fibre has size four. Instead consider the set
$\mathcal B=(z^Q)^P$ of block shadows neighbouring the base line.
There are four. The group $P$ is transitive on them, and the
intersection index $[P:P\cap Q]=4$ confirms that they are all its
incident blocks. Exactly $48$ of the $1920$ elements of $Q$ stabilize
$\mathcal B$ setwise. Put $T=G_{\mathcal B}$. We have $P\leq T$,
and the orbit of $z^Q$ under $T$ is all four members of
$\mathcal B$, since $P$ is already transitive. Also
$T_{z^Q}=T\cap Q$, by the block stabilizer equality. Hence
\[
                  |T|=4\cdot48=192=|P|.
\]
The four neighbouring blocks therefore determine each line.
Lemma~\ref{shadows} proves faithfulness of restriction of the full
geometry automorphism group to the point type.

The complete graph automorphism computation returns order
$88\,704\,000$ for each of these two graphs. Every returned generator
is checked against every adjacency. In GAP the group $D$ generated
by these permutations has this exact order, contains the actual
coset image $A$ of $G$, and satisfies $D'=A$. Completeness of the
graph computation is the explicitly retained nauty input; the
subgroup equality $D'=A$ is tested on the actual permutations.
Lemma~\ref{characteristic} now shows that any incidence isomorphism
between these geometries would conjugate their full specified $G$
actions. Such an isomorphism induces an automorphism of the abstract
group $G$. After composing with an element of $G$ to move the image
chamber back to the base chamber, it sends the unordered family of
parabolic stabilizers onto the other family. The acquired complete
$\Aut(G)$ comparison excludes this for representatives $14$ and $15$.
This proves that all $86$ representatives are distinct.

Finally, the complete panel inventory has exactly two two-transitive
actions in rank one, at subgroup anchors $586$ and $587$. Both
completed enumeration records have levels $(1,0)$, and both actions
have degree $176$, are faithful and are two-transitive. A rank one
geometry has no incidence between distinct types, so its isomorphism
type depends only on its number of objects. It has full automorphism
group $\operatorname{Sym}(176)$. The retained conjugators in the
acquired normalizer identify each point stabilizer with a parabolic
of a higher rank representative. They give explicit truncation
isomorphisms. There is consequently one rank one class, and none
which is not a truncation. Together with the acquired exhaustion in
higher ranks, this proves the theorem.
\end{proof}

\paragraph{Attribution and limits.}
The counts recover the classification asserted by Dimitri Leemans,
\emph{The residually weakly primitive geometries of HS},
Australasian Journal of Combinatorics 33 (2005), 231--236,
in the fixed local six-page edition. The complete subgroup catalogue,
the panel enumeration and the full automorphism action retain the
attributions and computational premises of their acquired reviews.
The shadow lemmas are elementary consequences of equivariance,
orbit--stabilizer and characteristic subgroups. Their application
above supplies the missing distinction between equivalence of group
actions and isomorphism of the geometries; it does not assert a new
classification or a solution to a new open question.

The exact retained computational inputs are GAP's permutation group,
coset action and subgroup algorithms and nauty's complete graph
automorphism algorithm. The supplied generators, adjacency tests,
stabilizer and derived subgroup equalities are additional finite
checks, not a substitute for nauty completeness. No Lean
formalization is claimed. The higher rank truncation counts, the
supplementary diagrams and full automorphism data, and the source's
two Hoffman--Singleton constructions remain separate obligations.
This component therefore gives no whole-manuscript completion or
compression measurement.
\end{document}
